\chapter{Completeness and the existence of integrals}
\label{chap:completeness}

In Chapter~\ref{chap:integrals} the weak integral of a curve was defined, but its existence was
always a hypothesis. In a Banach space this hypothesis is superfluous: continuous curves on a
compact interval have Riemann integrals. In a general locally convex space nothing of the sort
holds, and one has to pay for integrals with some completeness of $F$. The aim of this chapter
is to determine how much completeness is needed, and to draw the consequences for integrals
depending on parameters.

Two answers are classical. Weak integrals of continuous curves exist when $F$ is sequentially
complete, and weak integrals of $\Ck{1}$ curves exist under the weaker assumption of Mackey
completeness. The second one is the natural hypothesis for smooth calculus: it is what makes
parameter-dependent integrals of smooth maps smooth. Mackey completeness is also the completeness
notion of Kriegl and Michor \cite[\S 2]{KM97}; we stress that here it is only a hypothesis on
the model space, inside Michal--Bastiani calculus. Our main reference is \cite[\S1.1]{GN26}.

The chapter is organized as follows. In Section~\ref{sec:completeness-notions} we introduce
both completeness notions and compare them. In Section~\ref{sec:riemann} we study Riemann sums,
which in both cases provide the integral as a limit. Section~\ref{sec:existence} contains the
two existence theorems, Section~\ref{sec:ftc2} the second half of the fundamental theorem of
calculus, and Section~\ref{sec:param} the regularity of integrals depending on parameters.

\section{Completeness notions}
\label{sec:completeness-notions}

We state completeness through the topology of $F$, without fixing a uniform structure.

\begin{definition}[Sequential completeness]
  \label{def:seq-complete}
  \lean{LocallyConvex.Analysis.SeqComplete}
  \leanok
  A sequence $(u_n)$ in $F$ is Cauchy if for every neighborhood $W$ of $0$ there is $N$ with
  $u_n - u_m \in W$ for all $m, n \geq N$. The space $F$ is sequentially complete if every
  Cauchy sequence converges.
\end{definition}

\begin{definition}[Mackey completeness]
  \label{def:mackey-complete}
  \lean{LocallyConvex.Analysis.MackeyComplete}
  \leanok
  A sequence $(u_n)$ in $F$ is Mackey--Cauchy if there are a bounded set $B \subseteq F$ and
  real numbers $r_k \to 0$ with $u_l - u_k \in r_k B$ for all $k \leq l$. The space $F$ is
  Mackey complete if every Mackey--Cauchy sequence converges.
\end{definition}

Here bounded means bounded in the sense of von Neumann: absorbed by every neighborhood of $0$.
In other words, a Mackey--Cauchy sequence is Cauchy at a controlled rate, measured by a single
bounded set; in a normed space one may take for $B$ the unit ball, and the condition reads
$\|u_l - u_k\| \leq r_k$.

\begin{remark}
  Definition~\ref{def:mackey-complete} is one of the equivalent descriptions of Mackey
  completeness in \cite[1.1.18(e)]{GN26} and \cite[2.2, 2.14]{KM97}. Glöckner and Neeb take as
  definition the existence of weak integrals of Lipschitz curves; we prove in
  Chapter~\ref{chap:mackey-lipschitz} that both definitions agree, so the choice is a matter of
  convenience.
\end{remark}

\begin{proposition}[Comparison]
  \label{prop:completeness-chain}
  \lean{LocallyConvex.Analysis.SeqComplete.of_completeSpace, LocallyConvex.Analysis.SeqComplete.mackeyComplete}
  \leanok
  \uses{def:seq-complete, def:mackey-complete}
  Complete uniform additive groups are sequentially complete, and sequentially complete spaces
  are Mackey complete.
\end{proposition}

In particular, every complete locally convex space (for instance, a Fréchet or a Banach space)
is Mackey complete, and all the results of this chapter apply to it.

\begin{proof}
  \leanok
  The first claim amounts to comparing our notion of Cauchy sequence with the uniform one. For
  the second, it suffices to observe that a Mackey--Cauchy sequence is Cauchy: given a
  neighborhood $W$ of $0$, the bounded set $B$ is absorbed by $W$, so $r_k B \subseteq W$ for
  large $k$.
\end{proof}

We will need the following elementary fact, which is where local convexity enters.

\begin{lemma}[Convex hulls of bounded sets]
  \label{lem:conv-bounded}
  \lean{LocallyConvex.Analysis.isVonNBounded_convexHull}
  \leanok
  In a locally convex space, the convex hull of a bounded set is bounded.
\end{lemma}

\begin{proof}
  \leanok
  It is immediate that if $B$ is absorbed by a convex balanced neighborhood $W$ of $0$, so is
  $\conv B$; and these neighborhoods form a basis of neighborhoods of $0$.
\end{proof}

\section{Riemann sums}
\label{sec:riemann}

In both existence theorems the integral will be obtained as the limit of a sequence of Riemann
sums. We only need sums over uniform partitions.

\begin{definition}[Riemann sums]
  \label{def:riemann}
  \lean{LocallyConvex.Analysis.riemannSum}
  \leanok
  For $\gamma \colon \R \to F$ and $N \geq 1$, the Riemann sum of $\gamma$ over $[a, b]$ with
  $N$ equal subintervals and left endpoints is
  \[
    R_N(\gamma) = \frac{b-a}{N} \sum_{i=0}^{N-1} \gamma\Bigl(a + i\, \frac{b-a}{N}\Bigr).
  \]
\end{definition}

The basic estimate compares a Riemann sum with the sums of its refinements. It is stated with an
arbitrary convex set $V$, so that the same lemma serves both for neighborhoods of $0$ (the
sequentially complete case) and for multiples of a bounded set (the Mackey complete case).

\begin{lemma}[Refinements]
  \label{lem:riemann-refine}
  \lean{LocallyConvex.Analysis.riemannSum_mul_sub_mem, LocallyConvex.Analysis.tendsto_riemannSum_real, LocallyConvex.Analysis.map_riemannSum}
  \leanok
  \uses{def:riemann}
  Let $V \subseteq F$ be convex, and suppose that $\gamma(s) - \gamma(t) \in V$ whenever
  $s, t \in [a, b]$ and $0 \leq s - t \leq \frac{b - a}{N}$. Then
  $R_{Nm}(\gamma) - R_N(\gamma) \in (b - a)\, V$ for every $m \geq 1$. Moreover, for a continuous
  real function $g$ the Riemann sums $R_N(g)$ tend to $\int_a^b g$.
\end{lemma}

\begin{proof}
  \leanok
  The difference $R_{Nm} - R_N$ is $b - a$ times an average of increments of $\gamma$ over
  intervals of length at most $\frac{b - a}{N}$; since $V$ is convex, the average lies in $V$.
  The second claim is the classical convergence of Riemann sums of continuous functions.
\end{proof}

\begin{lemma}[Limits of Riemann sums]
  \label{lem:riemann-limit}
  \lean{LocallyConvex.Analysis.hasWeakIntegral_of_tendsto}
  \leanok
  \uses{lem:riemann-refine, def:weak-integral}
  If $\gamma$ is continuous on $[a, b]$ and the dyadic Riemann sums $R_{2^n}(\gamma)$ converge to
  $y$, then $y$ is the weak integral of $\gamma$ over $[a, b]$.
\end{lemma}

\begin{proof}
  \leanok
  For each continuous linear functional $\varphi$ we have
  $\varphi(R_{2^n}(\gamma)) = R_{2^n}(\varphi \circ \gamma)$, which converges both to
  $\varphi(y)$ and, by Lemma~\ref{lem:riemann-refine}, to $\int_a^b \varphi \circ \gamma$.
\end{proof}

\begin{remark}
  The dyadic sums $R_{2^n}(\gamma)$ form a \emph{sequence}, each term being a refinement of the
  previous one. This is why sequential notions of completeness suffice: no nets of partitions
  are needed.
\end{remark}

\section{Existence of integrals}
\label{sec:existence}

We can now state the two existence theorems. The construction is the same in both; only the
Cauchy estimate for the dyadic sums changes.

\begin{theorem}[Continuous curves, sequentially complete case]
  \label{thm:exists-seq}
  \lean{LocallyConvex.Analysis.exists_hasWeakIntegral, LocallyConvex.Analysis.exists_hasWeakIntegral_of_le}
  \leanok
  \uses{def:seq-complete, lem:riemann-limit}
  If $F$ is sequentially complete, every curve that is continuous on $[a, b]$ has a weak
  integral over $[a, b]$.
\end{theorem}

\begin{proof}
  \leanok
  By uniform continuity on the compact interval $[a, b]$, for each closed convex neighborhood
  $V$ of $0$ the increments of $\gamma$ over sufficiently short intervals lie in $V$. By
  Lemma~\ref{lem:riemann-refine}, the dyadic Riemann sums form a Cauchy sequence; its limit is
  the weak integral by Lemma~\ref{lem:riemann-limit}. This is \cite[1.1.15]{GN26}.
\end{proof}

\begin{theorem}[$\Ck{1}$ curves, Mackey complete case]
  \label{thm:exists-mackey}
  \lean{LocallyConvex.Analysis.exists_hasWeakIntegral_of_hasDirDerivAt_of_le, LocallyConvex.Analysis.ContDiffMBOn.exists_hasWeakIntegral, LocallyConvex.Analysis.sub_mem_smul_closure_convexHull}
  \leanok
  \uses{def:mackey-complete, lem:riemann-limit, lem:conv-bounded, lem:mean-value-inclusion}
  If $F$ is Mackey complete, every curve that is $\Ck{1}$ on a neighborhood of $[a, b]$ has a weak
  integral over $[a, b]$.
\end{theorem}

\begin{proof}
  \leanok
  Let $K$ be the closure of the convex hull of $\{s\,\gamma'(t) : s \in [0, 1],\ t \in [a, b]\}$.
  It is bounded by Lemma~\ref{lem:conv-bounded}, since $\gamma'([a, b])$ is compact. By the mean
  value inclusion (Lemma~\ref{lem:mean-value-inclusion}), $\gamma(s) - \gamma(t) \in h\, K$
  whenever $0 \leq s - t \leq h$. Then Lemma~\ref{lem:riemann-refine} gives
  \[
    R_{2^m}(\gamma) - R_{2^n}(\gamma) \in \frac{(b - a)^2}{2^n}\, K \qquad (n \leq m),
  \]
  so the dyadic Riemann sums form a Mackey--Cauchy sequence, with bounded set $K$ and
  $r_n = (b - a)^2 2^{-n}$. Its limit is the weak integral by Lemma~\ref{lem:riemann-limit}.
  Compare \cite[1.1.18]{GN26} and \cite[2.14]{KM97}.
\end{proof}

\begin{remark}
  In the continuous case the Cauchy estimate comes from uniform continuity, and it only says
  that the differences of dyadic sums are eventually small. In the $\Ck{1}$ case it comes from
  the mean value inclusion, and it says more: the differences lie in multiples of one fixed
  bounded set, with factors tending to $0$ at a known rate. This additional control is what
  allows us to replace sequential completeness by Mackey completeness.
\end{remark}

\section{The fundamental theorem of calculus, second part}
\label{sec:ftc2}

The first part of the fundamental theorem (Theorem~\ref{thm:ftc1}) recovers a $\Ck{1}$ curve
from the integral of its derivative. The second part, which we prove now, differentiates an
integral with respect to its upper limit.

\begin{theorem}[Fundamental theorem of calculus, II]
  \label{thm:ftc2}
  \lean{LocallyConvex.Analysis.hasDirDerivAt_weakIntegral_right_of_exists, LocallyConvex.Analysis.hasDirDerivAt_weakIntegral_right}
  \leanok
  \uses{lem:weak-integral-small, lem:weak-integral-algebra, prop:change-of-variables}
  Let $I \subseteq \R$ be an open interval, $a \in I$ and $\gamma$ continuous on $I$. If the weak
  integrals $\eta(s) = \int_a^s \gamma$ exist for $s \in I$, then $\eta$ is differentiable on
  $I$ with $\eta' = \gamma$. If $F$ is sequentially complete, these integrals always exist.
\end{theorem}

\begin{proof}
  \leanok
  By additivity in the interval and the affine change of variables $r \mapsto t + sr$,
  \[
    \frac{1}{s}\bigl(\eta(t + s) - \eta(t)\bigr) - \gamma(t)
      = \int_0^1 \bigl(\gamma(t + sr) - \gamma(t)\bigr)\, dr .
  \]
  Let $W$ be a closed convex neighborhood of $0$. By continuity of $\gamma$ at $t$, for small
  $s$ the integrand lies in $W$, and Lemma~\ref{lem:weak-integral-small} puts the integral in
  $2W$. The last claim follows from Theorem~\ref{thm:exists-seq}. This is \cite[1.1.12]{GN26}.
\end{proof}

\begin{remark}
  Note that the first assertion needs no completeness at all: what is assumed is the existence
  of the integrals $\eta(s)$, and the limit of the difference quotients is then controlled by
  Hahn--Banach through Lemma~\ref{lem:weak-integral-small}. Completeness of $F$ enters only to
  guarantee that existence.
\end{remark}

\section{Integrals depending on parameters}
\label{sec:param}

Let $U$ be an open subset of a locally convex space $E$, let $T \subseteq \R$ be open with
$[a, b] \subseteq T$, and let $f$ be a map on $U \times T$. We study the regularity of
$x \mapsto \int_a^b f(x, t)\, dt$. All the results follow from a single induction, in which the
existence of the integrals is isolated as a hypothesis.

\begin{theorem}[The induction behind parameter-dependent integrals]
  \label{thm:param-induction}
  \lean{LocallyConvex.Analysis.contDiffMBOn_weakIntegral_of_exists}
  \leanok
  \uses{thm:diff-under-integral, prop:param-continuity, thm:partials}
  Let $T \subseteq \R$ be open with $[a, b] \subseteq T$, and $m \in \N$. Suppose that, for
  every $\Ck{m}$ map $f$ on a product $U \times T$, the weak integrals $\int_a^b f(x, t)\, dt$
  exist. Then, for every $\Ck{k + m}$ map $f$ on $U \times T$, the map
  $x \mapsto \int_a^b f(x, t)\, dt$ is $\Ck{k}$ on $U$.
\end{theorem}

Here the hypothesis is understood for all open sets $U$ of all locally convex spaces $E$; this
uniformity is what makes the induction work.

\begin{proof}
  \leanok
  We argue by induction on $k$, for all locally convex spaces $E$ at once. The case $k = 0$ is
  Proposition~\ref{prop:param-continuity}. For the inductive step, differentiation under the
  integral sign (Theorem~\ref{thm:diff-under-integral}) gives
  \[
    \dd\Bigl(\int_a^b f(\cdot, t)\, dt\Bigr)(x)(h) = \int_a^b \dd f\bigl((x, t), (h, 0)\bigr)\, dt ,
  \]
  and the integrand is a $\Ck{k - 1 + m}$ map on the product $(U \times E) \times T$, to which
  the inductive hypothesis applies.
\end{proof}

\begin{remark}
  This single induction yields all the results on parameter-dependent integrals at once; the
  sources prove the sequentially complete case and the Mackey complete case separately.
\end{remark}

\begin{remark}[Lean]
  The quantification over all spaces $E$ is the reason for the universe restriction mentioned in
  Chapter~\ref{chap:differentiability}.
\end{remark}

Combined with the existence theorems of Section~\ref{sec:existence}, the induction gives the
following statements, with no existence hypotheses left.

\begin{corollary}[Parameter-dependent integrals]
  \label{cor:param}
  \lean{LocallyConvex.Analysis.ContDiffMBOn.contDiffMBOn_weakIntegral, LocallyConvex.Analysis.ContDiffMBOn.contDiffMBOn_weakIntegral_of_mackey, LocallyConvex.Analysis.SmoothMBOn.smoothMBOn_weakIntegral, LocallyConvex.Analysis.ContDiffMBOn.dMB_weakIntegral,
    LocallyConvex.Analysis.ContDiffMBOn.exists_hasWeakIntegral_slice,
    LocallyConvex.Analysis.ContDiffMBOn.dMB_weakIntegral_of_mackey}
  \leanok
  \uses{thm:param-induction, thm:exists-seq, thm:exists-mackey}
  Let $f$ be a map on $U \times T$ with $[a, b] \subseteq T$, and
  $I(x) = \int_a^b f(x, t)\, dt$.
  \begin{enumerate}
    \item If $F$ is sequentially complete and $f$ is $\Ck{k}$, then $I$ is $\Ck{k}$, and
    $\dd I(x)(h) = \int_a^b \dd f\bigl((x,t), (h,0)\bigr)\, dt$ for $k \geq 1$.
    \item If $F$ is Mackey complete and $f$ is $\Ck{k+1}$, then $I$ is $\Ck{k}$; if $f$ is
    $\Ck{2}$, the formula for $\dd I$ of (1) holds.
    \item If $F$ is Mackey complete and $f$ is smooth, then $I$ is smooth.
  \end{enumerate}
\end{corollary}

\begin{proof}
  \leanok
  Apply Theorem~\ref{thm:param-induction} with $m = 0$ and Theorem~\ref{thm:exists-seq} in the
  first case, and with $m = 1$ and Theorem~\ref{thm:exists-mackey} in the other two. This is
  \cite[1.3.23]{GN26}.
\end{proof}

\begin{remark}
  The loss of one derivative in the Mackey complete case has a simple origin: in that setting
  we only know that $\Ck{1}$ curves have integrals (Theorem~\ref{thm:exists-mackey}), so the
  induction has to be started with $m = 1$. For smooth maps nothing is lost, and this is why
  Mackey completeness is the natural hypothesis for smooth calculus.
\end{remark}
